\section{Méthode de recherche et d'analyse bibliographique}

\subsection{Périmètre et sources documentaires}

La démarche correspond à une revue narrative structurée avec recherche documentée et lecture critique. Elle privilégie les essais chez le poulet de chair et les travaux publiés entre 2018 et le 23 septembre 2026, tout en conservant des études antérieures directement pertinentes. Les articles sur la composition, l'extraction et la conservation des polyphénols sont recherchés séparément, sans restriction à l'espèce aviaire. Les études sur des molécules isolées et sur d'autres coproduits oléicoles sont identifiées comme données de contexte.

Deux requêtes reproductibles ont été exécutées dans PubMed par son interface E-utilities. La requête générale associe les termes \textit{olive leaf/leaves}, \textit{oleuropein} ou \textit{hydroxytyrosol}, recherchés dans les titres et résumés, aux termes \textit{broiler}, \textit{chicken} ou \textit{poultry}. Une requête complémentaire porte sur l'intestin, le microbiote et les contenus cæcaux. Le filtre de date de publication s'étend du 1\textsuperscript{er} janvier 2018 au 23 septembre 2026. Les chaînes exactes, traductions PubMed, identifiants et résultats bruts sont archivés dans le dossier documentaire.

La requête générale a retourné 48 notices, chacune examinée au niveau du titre : sept ont été classées comme candidates directement pertinentes, 28 comme candidates contextuelles et 13 comme hors périmètre. Douze titres comportent une ambiguïté nécessitant un contrôle du résumé ou du texte intégral. Ces catégories décrivent un dépistage initial ; elles ne constituent pas un nombre final d'études éligibles. La requête intestinale retourne neuf notices, déjà présentes dans les 48 notices générales : elles ne sont pas comptées comme neuf références supplémentaires.

Cette recherche est complétée par des requêtes ciblées sur le web, les pages des éditeurs, les références des documents locaux et les métadonnées Crossref. Deux appels Scite ont été utilisés : une interrogation groupée de onze DOI et une exploration des publications citant trois études prioritaires. Le graphe obtenu est limité à 45 liens et présente une couverture faible pour l'une des études ; il ne représente pas l'ensemble des citations disponibles. Les classifications automatiques des citations servent au repérage et ne remplacent pas l'examen du résultat original.

\subsection{Hiérarchie des études et extraction des résultats}

L'administration d'un extrait foliaire dans l'eau de boisson représente le niveau de proximité le plus élevé avec le sujet. Les extraits distribués dans l'aliment constituent une seconde catégorie. Les poudres de feuilles, l'oléuropéine ou l'hydroxytyrosol isolés, les grignons et les margines sont analysés séparément. Cette organisation réduit le risque d'attribuer à une infusion un effet observé avec un autre produit.

Pour chaque article retenu, l'extraction recherche, lorsque ces informations sont disponibles, la matrice, la préparation, la dose et son unité, la voie, les animaux et les répétitions, la durée, le contexte alimentaire ou thermique, les critères mesurés et la comparaison effectivement réalisée. Le résultat est rapproché des tableaux et des figures lorsqu'ils sont accessibles. Les résultats neutres ou défavorables sont conservés. Les incohérences entre valeurs de \(p\), lettres de comparaison et texte sont consignées ; elles ne sont pas corrigées par supposition.

Les recommandations ARRIVE 2.0 fournissent un cadre de lecture pour la transparence du plan expérimental, l'unité d'analyse, la randomisation, l'effectif et la présentation des résultats \citep{PercieduSert2020}. Elles ne constituent pas ici un score chiffré de qualité. En particulier, plusieurs oiseaux issus d'un même parquet ne sont pas nécessairement des répétitions indépendantes pour l'ingestion ou l'IC. Les prélèvements tissulaires, les répétitions analytiques et les unités auxquelles le traitement a été attribué doivent être distingués.

\subsection{Traçabilité et limites}

Le registre bibliographique rassemble 37 références sélectionnées, dont une référence méthodologique. Il distingue le texte intégral examiné, la lecture ciblée de sections, les extraits indexés et le résumé seul. Vingt PDF d'articles sont conservés dans un dossier séparé, et sept textes intégraux XML complètent les possibilités de vérification. Un téléchargement ne signifie pas qu'une lecture critique intégrale a été accomplie. Le fichier BibTeX est relié aux notices DOI ; les références non résolues par Crossref ont été contrôlées à partir du document original ou de PMC.

Cette étape ne couvre pas une interrogation exhaustive de CAB Abstracts, Scopus, Web of Science ou AGRIS. Les recherches ciblées de corrections n'ont pas trouvé de notice pertinente pour les quelques DOI contrôlés ; Scite n'a pas fourni de champ de notices éditoriales dans la réponse groupée. Ces constats ne garantissent donc pas l'absence de correction ou de rétractation dans toute la sélection. Le dépistage des notices, la vérification des textes encore manquants et la recherche de publications nouvelles devront être actualisés avant le dépôt en 2027.

Les résultats ne sont pas regroupés dans une méta-analyse : les matrices, unités de dose, durées, modèles et critères diffèrent, et plusieurs articles présentent des données ou indications statistiques insuffisantes pour une comparaison quantitative robuste. La synthèse cherche d'abord à identifier les convergences, les divergences et les conditions précises auxquelles chaque conclusion s'applique.
